\documentclass{article}
\usepackage{amsmath}
\usepackage{graphicx}
\usepackage{array}
\begin{document}
\begin{table}[h]
\caption{SVM com kernels Linear e Radial para base de dados 2 balanceada.}
\label{tab:sms_cls}
\begin{tabular}{c|c|c|c|c|c}
\hline
\textbf{Método} & \textbf{Kernel} & \textbf{$\gamma$} & \textbf{\begin{tabular}[c]{@{}c@{}}\emph{SNPs} selecionados\end{tabular}} & \textbf{\# \emph{SNPs} (V)} & \textbf{\begin{tabular}[c]{@{}c@{}}Acurácia\\média (\%)\end{tabular}} \\
\hline
SNPs causais & - & - & \begin{tabular}[c]{@{}c@{}}\textbf{1}, \textbf{2}, \textbf{3}, \textbf{4}, \textbf{5}\\\textbf{6}, \textbf{7}, \textbf{8}\end{tabular} & - & - \\
\hline
SMS & Linear & - & \begin{tabular}[c]{@{}c@{}}\textbf{1}, \textbf{2}, \textbf{3}, \textbf{5}, \textbf{8}\\19, 37, 59, 71\end{tabular} & 9(5) & 76.50 \\
\hline
SMS & Radial & 0.001 & \begin{tabular}[c]{@{}c@{}}\textbf{1}, \textbf{2}, \textbf{3}, \textbf{4}, \textbf{5}\\\textbf{6}, \textbf{7}, \textbf{8}, 11, 16\\20, 24, 27, 31, 32\\43, 53, 54, 59, 64\\76, 78, 82, 87, 89\\92, 97, 99\end{tabular} & 28(8) & 76.40 \\
\hline
SMS & Radial & 0.01 & \begin{tabular}[c]{@{}c@{}}\textbf{1}, \textbf{2}, \textbf{3}, \textbf{4}, \textbf{5}\\\textbf{6}, \textbf{8}, 19, 24, 27\\31, 38, 43, 54, 71\\79, 89, 92, 97, 99\end{tabular} & 20(7) & 78.50 \\
\hline
SMS & Radial & 0.1 & \begin{tabular}[c]{@{}c@{}}\textbf{1}, \textbf{2}, \textbf{3}, \textbf{4}, \textbf{5}\\\textbf{8}, 11, 16, 27, 37\\54, 59, 89, 92, 99\end{tabular} & 15(6) & 80.90 \\
\hline
SMS & Radial & 1.0 & \begin{tabular}[c]{@{}c@{}}\textbf{1}, \textbf{2}, \textbf{3}, \textbf{5}, 27\\59\end{tabular} & 6(4) & 77.60 \\
\hline
União & - & - & \begin{tabular}[c]{@{}c@{}}\textbf{1}, \textbf{2}, \textbf{3}, \textbf{4}, \textbf{5}\\\textbf{6}, \textbf{7}, \textbf{8}, 11, 16\\19, 20, 24, 27, 31\\32, 37, 38, 43, 53\\54, 59, 64, 71, 76\\78, 79, 82, 87, 89\\92, 97, 99\end{tabular} & 33(8) & - \\
\hline
Interseção & - & - & \begin{tabular}[c]{@{}c@{}}\textbf{1}, \textbf{2}, \textbf{3}, \textbf{5}\end{tabular} & 4(4) & - \\
\hline
Valor-p bruto & - & - & \begin{tabular}[c]{@{}c@{}}\textbf{1}, \textbf{2}, \textbf{3}, \textbf{4}, \textbf{5}\\\textbf{8}, 11, 42, 43, 74\\89\end{tabular} & 11(6) & - \\
\hline
Valor-p corrigido & - & - & \begin{tabular}[c]{@{}c@{}}\textbf{1}, \textbf{2}, \textbf{3}, 11\end{tabular} & 4(3) & - \\
\hline
\end{tabular}
\end{table}
\end{document}
